\documentclass[xcolor=dvipsnames,aspectratio=169,notes]{beamer}

\title{Reviving the Media Station:\\A Study in Reverse Engineering}
\author{Nathanael Gentry}
\institute{School of Engineering}
\date{April 13, 2021}

\usepackage[T1]{fontenc}
\usepackage{caption}
\renewcommand{\figurename}{}

% Unused ultimately, but still interesting:
% https://tex.stackexchange.com/questions/569923/conflict-of-resizebox-and-verbatim-mode
\newsavebox{\tablebox}

% ~~~~~ LIGHT THEME ~~~~~~~~~~~~~~~~~~~~~~~~~~~~~~~
% \usetheme{Berkeley}
% \usecolortheme{whale}

% \definecolor{LUblue}{rgb}{0.03921568627, 0.14509803921, 0.30588235294} % UBC Blue (primary)
% \definecolor{LUred}{rgb}{0.6,0,0} % UBC Grey (secondary)

% \setbeamercolor{palette primary}{bg=LUblue,fg=white}
% \setbeamercolor{palette secondary}{bg=LUblue,fg=white}
% \setbeamercolor{palette tertiary}{bg=LUblue,fg=white}
% \setbeamercolor{palette quaternary}{bg=LUblue,fg=white}
% \setbeamercolor{structure}{fg=LUblue} % itemize, enumerate, etc
% \setbeamercolor{section in toc}{fg=LUblue} % TOC sections

% \usebackgroundtemplate{\includegraphics[width=\paperwidth,height=\paperheight]{solidwhite.jpg}}

% ~~~~~ DARK THEME ~~~~~~~~~~~~~~~~~~~~~~~~~~~~~~~
% https://tex.stackexchange.com/questions/57477/beamer-dark-theme
% \usetheme{Warsaw}

\setbeamercolor{normal text}{fg=white,bg=black!90}
\setbeamercolor{structure}{fg=white}

\setbeamercolor{alerted text}{fg=red!85!black}

\setbeamercolor{item projected}{use=item,fg=black,bg=item.fg!35}

\setbeamercolor*{palette primary}{use=structure,fg=structure.fg}
\setbeamercolor*{palette secondary}{use=structure,fg=structure.fg!95!black}
\setbeamercolor*{palette tertiary}{use=structure,fg=structure.fg!90!black}
\setbeamercolor*{palette quaternary}{use=structure,fg=structure.fg!95!black,bg=black!80}

\setbeamercolor*{framesubtitle}{fg=white}

\setbeamercolor*{block title}{parent=structure,bg=black!60}
\setbeamercolor*{block body}{fg=black,bg=black!10}
\setbeamercolor*{block title alerted}{parent=alerted text,bg=black!15}
\setbeamercolor*{block title example}{parent=example text,bg=black!15}

\usebackgroundtemplate{\includegraphics[width=\paperwidth,height=\paperheight]{demossblue.jpg}}

% https://tex.stackexchange.com/questions/411313/using-listings-package-with-beamer-block-environments-for-displaying-source-code
\usepackage{listings}
\lstset{numbers=left, numberstyle=\tiny, stepnumber=1,firstnumber=1,
    numbersep=5pt,language=C,
    stringstyle=\ttfamily,
    basicstyle=\footnotesize,
    showstringspaces=false
}

% https://tex.stackexchange.com/questions/48473/best-way-to-give-sources-of-images-used-in-a-beamer-presentation
\usepackage[absolute,overlay]{textpos}

\setbeamercolor{framesource}{fg=gray}
\setbeamerfont{framesource}{size=\footnotesize}

\newcommand{\source}[1]{\begin{textblock*}{4cm}(11.5cm,8cm)
    \begin{beamercolorbox}[ht=0.5cm,right]{framesource}
        \usebeamerfont{framesource}\usebeamercolor[fg]{framesource} Source: {#1}
    \end{beamercolorbox}
\end{textblock*}}

% https://tex.stackexchange.com/questions/16793/global-setting-of-spacing-between-items-in-itemize-environment-for-beamer
\let\olditem\item
\renewcommand{\item}{\olditem\vspace{8pt}}

% https://tex.stackexchange.com/questions/262115/bold-serif-headings-in-beamer
\setbeamerfont{frametitle}{series=\bfseries,parent=structure}

% https://www.patrickbaylis.com/blog/2018-10-11-beamer-resizing/
\usepackage{adjustbox}
\makeatletter
\newcommand{\fitimage}[2][\@nil]{
  \begin{figure}
    \begin{adjustbox}{width=0.9\textwidth, totalheight=\textheight-2\baselineskip-2\baselineskip,keepaspectratio}
      \includegraphics{#2}
    \end{adjustbox}
    \def\tmp{#1}%
    \ifx\tmp\@nnil
    \else
    \caption{#1}
    \fi
  \end{figure}
}
\makeatother

% \usepackage{listings}
\usepackage{adjustbox}

\begin{document}

\frame{\titlepage}

\usebackgroundtemplate{\includegraphics[width=\paperwidth,height=\paperheight]{solidbluenolu.jpg}}
\frame{}

\begin{frame}
  \fitimage{gsocsun}
  \source{\cite{GoogleSummerCode2020}}
\end{frame}

\begin{frame}
  \fitimage{gsoc}
  \source{\cite{GoogleSummerCode2020}}
\end{frame}

\begin{frame}
  \fitimage{scummvm}
  \source{\cite{ScummVM}}
\end{frame}

\begin{frame}
  \vspace{-2em}
  \begin{figure}
    \includegraphics[scale=0.4]{scummvm}
  \end{figure}
  \vspace{1em}\pause
  \begin{itemize}
    \item Open-source reimplementation of game interpreters\pause
    \item Original game data playable on newer and diverse platforms
  \end{itemize}
  \source{\cite{ScummVM}}
\end{frame}

\begin{frame}
  \fitimage{d4cover}
  \source{\cite{TipsTricksDirector1994}}
\end{frame}

\begin{frame}
\fitimage{candle}
\end{frame}

\begin{frame}
  \fitimage{alien}
\end{frame}

\newcommand{\github}{\begin{frame}
  \fitimage{seuss2}
  \vspace{-1.5em}
  \begin{center}
    \huge \url{https://npjg.github.io}
  \end{center}
\end{frame}}
\github

\begin{frame}
  \fitimage{asblogo}
  \source{\cite{DisneyAnimatedStoryBook2020}}
\end{frame}

\begin{frame}
  \fitimage{msitrans}
  \source{\cite{flurryDisneyAnimatedStoryBook1994}}
\end{frame}

\begin{frame}
\begin{center}
  ``I lost my copy of the source many years ago when a backup drive failed. I was
  hoping to do a similar thing to what ScummVM does...''
\end{center}
\end{frame}

\begin{frame}{Research Question}\pause
  \begin{center}
    \huge
    How can an interactive multimedia format be reverse-engineered from scratch?
  \end{center}
\end{frame}

% To show how I have begun answering this question...

\begin{frame}{Counting}\pause
\begin{itemize}
  \item \textbf{Base-10} (decimal) uses 10 digits --\pause{} the numbers 0-9.\\\pause
  \item \textbf{Base-2} (binary) uses 2 digits --\pause{} the bits \texttt{0} and \texttt{1}.\\\pause
  \item \textbf{Base-16} (hexadecimal) uses 16 digits --\pause{} the numbers 0-9\pause{} and
    letters A-F.
\end{itemize}
\end{frame}

\begin{frame}{Bytes}\pause
  \begin{center}
    \large

    A \textbf{byte} is a collection of 8 bits.\\\pause
    Each byte represents a value from $0 = (2^0 - 1)$\pause{} to $255 = (2^8 - 1)$.\pause

    \vspace{1em}
    For example,\pause
    \begin{equation*}
      (10101010)_2\pause \equiv \mathbf{(AA)}_{16}\pause \equiv (170)_{10}.\pause
    \end{equation*}

    Additionally,\pause
    \begin{equation*}
      \mathbf{(AA)_{16}}\pause\equiv  \textbf{\texttt{0xAA}}\pause \equiv  \textbf{\texttt{0xaa}}.
    \end{equation*}
  \end{center}
\end{frame}

\begin{frame}{Files}\pause
  \begin{center}
    A \textbf{file} is a collection of bytes united under one name.\\\pause
    Each byte has a unique \textbf{address} -- its position in the file.\pause
  \end{center}
\end{frame}


\begin{frame}
\begin{itemize}
\item Look for patterns
\item Let nothing go unnoticed
\item Write throwaway code
\end{itemize}
\end{frame}

\begin{frame}{Visit Me}
  GitHub: \url{https://github.com/npjg}
\end{frame}

\begin{frame}[allowframebreaks]{References}
  \bibliographystyle{IEEEtran}
  \bibliography{MediaStation}
\end{frame}
\end{document}
